% International Journal of Fracture submission source
% Generated from the scientifically frozen English master on 2026-09-25.
% Springer Nature template: December 2024; IJF requests the iicol option.
\documentclass[pdflatex,iicol,sn-mathphys-ay]{sn-jnl}

\usepackage{graphicx}
\usepackage{tikz}
\usetikzlibrary{arrows.meta,calc}
\usepackage{amsmath,amssymb,mathtools,bm}
\usepackage{booktabs,array}
\usepackage{xcolor}

\newcommand{\dd}{\mathop{}\!\mathrm{d}}
\newcommand{\ii}{\mathrm{i}}
\newcommand{\ee}{\mathrm{e}}
\newcommand{\sigmap}{\sigma^{\prime}}
\newcommand{\deltap}{\delta^{\prime}}
\newcommand{\Jp}{J^{\prime}}
\newcommand{\jump}[1]{\left[\!\left[#1\right]\!\right]}
\newcommand{\resultfigure}[4][0.98\linewidth]{%
  \begin{figure}[htbp]
    \centering
    \includegraphics[width=#1]{#2}%
    \caption{#3}
    \label{#4}
  \end{figure}}

\begin{document}

\title[Branching initiation from a V-shaped interfacial shear crack]{Analytical Model for Branching Initiation from the Corner of a V-Shaped Interfacial Shear Crack}

\author[1]{\fnm{A. O.} \sur{Kaminsky}}
\author[1]{\fnm{M. F.} \sur{Selivanov}}
\author[1,2]{\fnm{M. V.} \sur{Dudyk}}
\author[2]{\fnm{T. V.} \sur{Polishchuk}}

% TODO before final submission: identify the corresponding author with \author* and add \email{...}.
\affil[1]{\orgname{S. P. Timoshenko Institute of Mechanics, National Academy of Sciences of Ukraine}, \orgaddress{\city{Kyiv}, \country{Ukraine}}}
\affil[2]{\orgname{Uman National University}, \orgaddress{\city{Uman}, \country{Ukraine}}}

\abstract{Branching initiation from the corner of a stationary V-shaped interfacial shear crack at a kinked interface between two isotropic elastic materials is studied under plane strain. The crack consists of two symmetric branches whose faces remain in frictionless contact. Its corner singular field is used as the outer field for a small tensile process zone directed along the bisector of one material sector. Physical admissibility requires simultaneous compressive contact on the interfacial crack and a tensile normal field in the material containing the zone. Mellin transformation, asymptotic matching, and the Wiener--Hopf method yield closed analytical relations for the equilibrium zone length, opening, work of separation, and a zone-related energy parameter. Their dependence on interface angle, elastic contrast, and loading is examined. The process zone weakens, but does not eliminate, the corner singularity. For fixed geometry, the deformation and energy measures are uniquely related and provide equivalent branching-initiation criteria when their critical parameters are calibrated consistently. Once a critical material characteristic is specified, the model determines the corresponding critical external load.}

\keywords{piecewise homogeneous body, kinked interface, V-shaped interfacial shear crack, process zone, branching initiation, Wiener--Hopf method}

\maketitle

\section{Introduction}
\label{sec:introduction}

Corners and kinks in interfaces between dissimilar materials occur in many
adhesive and multimaterial joints. In such configurations, the interface
geometry itself can create local stress concentrations that strongly affect
damage initiation~\citep{GacoinEtAl2009}. In an idealized elastic setting, a
bimaterial interface corner is characterized by a power-law eigenfield whose
exponent and angular structure depend on the geometry and elastic contrast.
Classical existence conditions and characteristics of such singularities were
established for bonded isotropic wedges~\citep{BogyWang1971,DempseySinclair1981};
subsequent studies showed that changing the angle can substantially alter both
the singularity exponent and the intensity of the local field
\citep{MohammedLiechti2001}. Methods for determining singular-field parameters
in more complex three-dimensional configurations continue to be developed
\citep{MiyazakiFukudaNoda2025}, while small-scale relations between the
parameters of an undamaged corner field and those of a short crack emanating
from an interface corner have also been proposed
\citep{LyuZhangMai2026}. Thus, fracture behavior depends not only on the corner singularity itself,
but also on its interaction with a local nonlinear or process zone.

For initially undamaged bimaterial corners, the transition from an elastic
singularity to local fracture has been studied both through parameters of the
asymptotic field and through models with a finite cohesive or process zone.
\citet{MohammedLiechti2000} combined experimental measurements with cohesive-zone
modeling of crack nucleation for different bimaterial-joint angles, whereas
\citet{LabossiereDunn2001} examined fracture initiation at three-dimensional
bimaterial corners. \citet{AkisanyaMeng2003} showed that the use of a
critical parameter of the elastic corner field requires sufficient scale
separation between the nonlinear zone and the region dominated by the
asymptotic field. In modern cohesive formulations,
the same requirement is expressed through the ratio of a characteristic
fracture-zone length to the geometric scale of the problem
\citep{GormanThouless2022,ThoulessGoutianos2023}. For local analytical models, this requirement motivates the small-scale
assumption that the process zone be small relative to the scale of the outer
singular field to which it is matched.

A second group of studies concerns existing interface cracks and their
possible departure from the interface. For an ideally brittle description,
classical energetic analyses established conditions for initial crack kinking
out of an interface~\citep{HeHutchinson1989}; corresponding relations for
stress intensity factors and energy release rates were obtained for
anisotropic bimaterials~\citep{Wang1994}. Cohesive models show that path
selection may depend not only on fracture toughness, but also on cohesive
strength and the characteristic fracture-zone scale
\citep{ParmigianiThouless2006}. Analytical models of small process zones near
interface cracks provide the zone length, opening, and energy characteristics
\citep{KaminskyEtAl2023}; accounting for the
nonsingular $T$-stress further shows that higher-order terms of the outer
field can produce appreciable quantitative corrections
\citep{KaminskyDudykChornoivan2026}. These results provide the methodological basis for the present process-zone
approach, although the geometry considered here is different.

For corners of kinked interfaces, analytical models of process zones and
cracks have been developed in several closely related settings.
\citet{DudykDikhtyarenko2012} studied the initial stage of an interface crack
kinking from a corner point, whereas \citet{KaminskyKipnisPolishchuk2012}
considered an initial fracture process zone arising directly at the corner of
a V-shaped interface. For an interface crack with contacting
faces emanating from a corner point of a polygonal interface, parameters of a
small fracture zone were determined in
\citep{KaminskyDudykFenkiv2023}. A related model addresses crack initiation near a kinked interface in the
presence of a plastic connecting layer~\citep{KaminskyDudykPolishchuk2024}. Thus, the analytical framework for local process zones near interface
corners and interface cracks is well developed. The present study applies
this framework to the specific singular field generated by a stationary
contacting interfacial shear crack at the corner of a V-shaped interface.

Such an outer field was obtained by \citet{NazarenkoKipnis2022Stress} for a
symmetric interfacial shear crack with faces in frictionless contact. In the present work, the two symmetric
branches of this defect are treated as a single V-shaped interfacial shear
crack. Each branch has length $l$, so the total length of the two branches is
$2l$. In the cracked configuration, a power-law singular field remains near
the corner; its exponent and amplitude depend on the geometry, elastic
contrast, and $l$. This field provides the immediate outer asymptotics for
the local process-zone problem considered below.

For this configuration, \citet{NazarenkoKipnis2022Equilibrium} also determined
the stress intensity factor for in-plane shear (Mode~II) at the tips of the
V-shaped crack branches, formulated the limiting-equilibrium condition, and
showed that this equilibrium is unstable with respect to further interfacial
propagation. That mechanism is localized
at the crack-branch tips and is potentially competitive with the mechanism studied
here at the corner. In the present work, the V-shaped interfacial crack is
treated as a prescribed stationary configuration, and the applied loading is
assumed insufficient to cause loss of its local stability. Under these
conditions, the singular corner field can support a small tensile process
zone in one of the materials. The corner and crack-branch-tip mechanisms are
therefore treated as competing mechanisms rather than as sequential stages
of a prescribed fracture scenario.

This work focuses on the local process zone near the corner of a V-shaped
interfacial shear crack and on using its characteristics to predict
branching initiation into one of the materials. Under the small-scale
approximation, the zone length in material $i$ ($i=1,2$) is denoted by $d_i$
and is assumed to be much smaller than the crack-branch length $l$. The crack field in the absence of the zone is used as the outer field, and
the leading-order feedback of the zone on the crack-branch tips is neglected.
The first objective is to determine which material can contain a physically
admissible zone. We then determine its equilibrium length $d_i$, opening
$\delta_i$, work $W_i$ done against resistance to separation, zone-related
energy parameter $J_i=\dd W_i/\dd d_i$, and the corresponding change in the
local corner singularity.

Accordingly, we construct an analytical model for the quasi-static family of
local equilibrium states of the process zone and use it to formulate
conditions for branching initiation of the V-shaped interfacial shear crack
at its corner. Under plane strain and symmetric loading, a small tensile
process zone is considered in one of the materials and is oriented along the
bisector of the corresponding angular sector. The zone is modeled as a line
of normal displacement discontinuity subjected to a constant normal cohesive
traction $\sigma_i$. Mellin transformation, asymptotic matching, and the Wiener--Hopf method are
used to determine the physical-admissibility conditions, equilibrium zone
length and opening, work done against resistance to separation, zone-related
energy parameter, and exponent of the residual corner singularity. The opening
and energy parameter may then be used in deformation and energy criteria,
respectively, for initiation of the branched crack once the corresponding
critical characteristic is specified independently. Subsequent propagation and the finite branching path
are outside the scope of this local formulation.

The remainder of the paper is organized as follows. Section~\ref{sec:problem}
formulates the local boundary-value problem and specifies the geometry,
boundary conditions, and asymptotic fields for the V-shaped interfacial crack
and the process zone. Section~\ref{sec:derivation} uses the Mellin transform,
asymptotic matching, and the Wiener--Hopf method to determine the equilibrium
zone length, opening, and energy characteristics. Section~\ref{sec:numerical-results}
examines the physically admissible regimes, the dependence of the zone
parameters on geometry and loading, and the changes in the corner singularity
exponents. Section~\ref{sec:discussion} discusses branching-initiation
criteria, competing fracture mechanisms, the role of the residual
singularity, the range of applicability, and possible generalizations of the
model. The main results are summarized in Section~\ref{sec:conclusions}.


\section{Problem formulation}
\label{sec:problem}
\providecommand{\AsymptoticsAppendix}{Appendix~\ref{app:asymptotics}}

\subsection{Geometry and initial configuration}
\label{subsec:geometry}

Consider a piecewise homogeneous body under plane strain composed of two
dissimilar isotropic elastic materials joined along a kinked interface
(Figure~\ref{fig:failure-stages}).
Material $i$ has Young's modulus $E_i$ and Poisson's ratio $\nu_i$, where
$i=1,2$ identifies the material. The angle between the two interface rays,
measured through material~2, is $2\alpha$, where $\alpha$ is the half-angle and
\begin{equation*}
  0<\alpha<\pi, \qquad \alpha\ne\frac{\pi}{2}.
\end{equation*}
The value $\alpha=\pi/2$ corresponds to a straight interface and is treated
only as a limiting case.

The origin of the polar coordinate system $(r,\theta)$ is placed at the
interface corner $O$, where $r$ is the distance from $O$ and $\theta$ is
measured from the bisector of the sector occupied by material~2. Symmetry
about this bisector allows the analysis to be restricted to the half-plane
$0\leq\theta\leq\pi$. Material~2 occupies $0<\theta<\alpha$, material~1
occupies $\alpha<\theta<\pi$, and the interface lies along
$\theta=\alpha$. The radial and circumferential displacement components
are denoted by $u_r$ and $u_\theta$, respectively. The circumferential normal
stress is $\sigma_\theta$, and the shear stress in polar coordinates is
$\tau_{r\theta}$. On a radial line, these are the normal and tangential
traction components, respectively.

The loading produces a local field that is symmetric about the bisector.
For the intact interface, the singular part of the stress field along the
interface is
\begin{equation*}
\begin{aligned}
  \tau_{r\theta}(r,\alpha)
    &= Cg_1 r^{\lambda_0}+o\!\left(r^{\lambda_0}\right),\\
  \sigma_\theta(r,\alpha)
    &= Cg_2 r^{\lambda_0}+o\!\left(r^{\lambda_0}\right),
    \qquad r\to0.
\end{aligned}
\end{equation*}
Here, $C$ is the amplitude of the local singular field and incorporates the
effects of the applied loading and the global body geometry; $g_1$ and
$g_2$ are angular coefficients that depend on $\alpha$ and the elastic
properties of the materials; and $\lambda_0\in(-1,0)$ is the singularity
exponent of the intact corner. The notation $o(r^{\lambda_0})$ denotes a
remainder whose ratio to $r^{\lambda_0}$ tends to zero as $r\to0$.
The exponent $\lambda_0$ is the smallest real root of
$D_0(-1-\lambda_0)=0$ in the open interval $(-1,0)$ that is continuously
connected to the physical solution branch. Here, $D_0$ is the characteristic
determinant for the intact corner. Expressions for $D_0$, $g_1$, and $g_2$
are given in \AsymptoticsAppendix.

The initial configuration considered here contains an already existing
V-shaped interfacial shear crack formed by two symmetric branches emanating
from the corner $O$ along the kinked interface. The crack faces are in
frictionless contact. Each branch has length $l$, so the total length of the
two branches is $2l$; below, $l$ is referred to as the crack half-length
(the branch length). The contact state is consistent with
\begin{equation*}
  Cg_2<0,
\end{equation*}
for which the normal stress along the interface is compressive. The
V-shaped interfacial crack is treated as a prescribed stationary
configuration and $l$ as a prescribed fixed parameter. The applied loading
is assumed insufficient to cause loss of local stability at the outer crack
tips. Interfacial propagation from those tips is outside the local problem
near $O$ and is considered below only as a competing fracture mechanism.

The prescribed interfacial crack retains a power-law stress singularity at
the corner~\citep{NazarenkoKipnis2022Stress}. This field can support a small
tensile process zone in one of the materials; its characteristics are used
below to formulate the conditions for branching initiation at the corner. Under the assumed geometric and loading symmetry, the bisector of the sector
occupied by that material is a natural direction for the zone. Along this
direction, the local Mode~II (in-plane shear) component at the zone tip
vanishes by symmetry, consistent with the principle of local
symmetry~\citep{BarenblattCherepanov1961,GoldsteinSalganik1974}. Zones with
arbitrary orientations are not considered. The bisector direction is
therefore part of the assumed symmetric model and is not obtained by
separately optimizing the branching direction.

Let $\beta_i$ denote the polar angle of the bisector in material~$i$:
\begin{equation*}
  \beta_1=\pi,\qquad \beta_2=0.
\end{equation*}
The subscript $i$ in $\beta_i$, $d_i$, $\sigma_i$, and $Q_i$ identifies
the material throughout.

The process zone in material~$i$ is modeled as a normal displacement
discontinuity of length $d_i$ subjected to a constant normal cohesive
traction $\sigma_i>0$, which represents the material's resistance to
separation. Figure~\ref{fig:failure-stages} shows the zone in material~1
as a solid line and the alternative zone in material~2 as a dashed line.
These are alternative configurations, considered separately. For each
admissible direction, the solution defines a quasi-static family of
equilibrium process-zone states; the zone opening and energy characteristic
are then used to formulate branching-initiation criteria.

\begin{figure}[htbp]
  \centering
  \includegraphics[width=\linewidth]{Fig1.pdf}
  \caption{Local configurations near the interface corner:
  (a) intact kinked interface; (b) prescribed stationary V-shaped
  interfacial shear crack with two branches of length $l$ each (total branch
  length $2l$); (c) process zone of length $d_1$
  in material~1 and the alternative zone of length $d_2$ in material~2
  (dashed line). The panels compare the corresponding local configurations
  and do not prescribe the order in which they occur.}
  \label{fig:failure-stages}
\end{figure}

\subsection{Local problem and boundary conditions}
\label{subsec:boundary-conditions}

Let $L$ be a characteristic outer dimension of the body. We assume the
separation of scales
\begin{equation*}
  d_i\ll l\ll L
\end{equation*}
and neglect friction between the interfacial crack faces. Within this
quasi-static formulation, the outer-field amplitude $C$ may vary, whereas
the crack half-length $l$, and hence the total length $2l$ of the two
branches, remain fixed. Changes in loading alter the equilibrium zone length
$d_i$ and the associated local characteristics. In the local problem, the
piecewise homogeneous body is replaced by an infinite plane, and the two
branches of the V-shaped interfacial crack are treated as semi-infinite.
Their finite length and the applied loading enter through an asymptotic
matching condition~\citep{Nayfeh1973Perturbation}. This is a one-way local
approximation: the known field of the stationary interfacial crack loads the
inner problem near $O$, whereas the feedback of the small process zone on
the fields at the crack-branch tips is not determined.

For a field quantity $f$ defined on both sides of the ray
$\theta=\alpha$, its jump across the interface is defined by
\begin{equation*}
  \jump{f}
  =f(r,\alpha+0)-f(r,\alpha-0),
\end{equation*}
where $\alpha+0$ and $\alpha-0$ denote the limiting values approached
from materials~1 and~2, respectively. Symmetry, traction continuity, and
frictionless contact give
\begin{equation}
\label{eq:interface-boundary-conditions}
\begin{aligned}
  &\theta\in\{0,\pi\}: &&\tau_{r\theta}=0,\\
  &\theta=\alpha: &&
    \jump{\sigma_\theta}=0,
    \quad \jump{\tau_{r\theta}}=0,
    \quad \tau_{r\theta}=0,
    \quad \jump{u_\theta}=0.
\end{aligned}
\end{equation}
The condition $\jump{u_\theta}=0$ imposes normal contact between the crack
faces, whereas $\tau_{r\theta}=0$ expresses the absence of friction.

On the bisector of the material sector containing the process zone, the
following mixed boundary conditions apply:
\begin{equation}
\label{eq:process-zone-boundary-conditions}
\begin{aligned}
  &\theta=\beta_i,\quad 0<r<d_i:
    &&\sigma_\theta(r,\beta_i)=\sigma_i,\\
  &\theta=\beta_i,\quad r>d_i:
    &&u_\theta(r,\beta_i)=0,\\
  &\theta=\beta_j,\quad r>0,\quad j\ne i:
    &&u_\theta(r,\beta_j)=0.
\end{aligned}
\end{equation}
Here, $j=3-i$ identifies the other material. Thus, the symmetry
condition on $u_\theta$ is replaced by a prescribed constant normal
traction only over $0<r<d_i$ along the active direction
$\theta=\beta_i$.

At intermediate distances $d_i\ll r\ll l$, the solution must match the
asymptotic field of the interfacial crack in the absence of a process
zone~\citep{NazarenkoKipnis2022Stress}. In the inner problem, this matching
condition takes the form
\begin{equation}
\label{eq:matched-far-field}
  \sigma_\theta(r,\beta_i)
    =CQ_i r^\lambda+o\!\left(r^{-1}\right),
    \qquad r\to\infty.
\end{equation}
Here, $Q_i$ is the amplitude coefficient for material~$i$, which depends
on the interface angle, the elastic properties, and the branch length
$l$, as specified in \AsymptoticsAppendix. The exponent
$\lambda\in(-1,0)$ characterizes the corner singularity in the
configuration with the prescribed stationary V-shaped interfacial crack. It is the smallest real root of
$D(-1-\lambda)=0$ in $(-1,0)$ that is continuously connected to the
physical branch, where $D$ is the characteristic determinant for the
cracked configuration. Expressions for $D$ and $Q_i$ are given in
\AsymptoticsAppendix. The limit $r\to\infty$ is understood in the
inner-scale sense of asymptotic matching, rather than as the far field
of the finite body; the remainder $o(r^{-1})$ decays faster than $r^{-1}$
in that limit.

For the subsequent Wiener--Hopf construction, we use superposition as in
the related problem~\citep{KaminskyDudykPolishchuk2024}. We subtract the known field of the stationary V-shaped interfacial crack
without a process zone and introduce the correction normal stress
\begin{equation*}
  \widehat{\sigma}_i(r)
  =\sigma_\theta(r,\beta_i)-CQ_i r^\lambda.
\end{equation*}
The matching condition then gives
\begin{equation*}
  \widehat{\sigma}_i(r)=o(r^{-1}),
  \qquad r\to\infty,
\end{equation*}
while on the process-zone segment $0<r<d_i$,
\begin{equation*}
  \widehat{\sigma}_i(r)
  =\sigma_i-CQ_i r^\lambda.
\end{equation*}
The field of the stationary V-shaped crack without a process zone satisfies
the bisector symmetry condition $u_\theta=0$, so subtracting it does not alter the displacement-gradient
transform used below. The functional problem is therefore formulated for the correction field,
with the known power-law term included in the prescribed part.

A process zone can form in material~$i$ only if the normal stress on its
bisector is tensile. Equation~\eqref{eq:matched-far-field} therefore gives
the physical admissibility condition
\begin{equation*}
  CQ_i>0.
\end{equation*}
The model thus requires both $Cg_2<0$, for contact between the
interfacial crack faces, and $CQ_i>0$, for formation of a process zone
in material~$i$. For the parameter sets examined below, the numerical
analysis shows that these conditions can hold simultaneously for only one
of the two materials; only that material admits a physically admissible
process zone.

\subsection{Asymptotic fields near the process-zone tip}
\label{subsec:zone-tip-asymptotics}

Near the outer tip of the displacement discontinuity, the stress and normal
displacement gradient have the standard inverse-square-root asymptotics for a
homogeneous material:
\begin{equation}
\label{eq:zone-tip-asymptotics}
\begin{aligned}
  &\theta=\beta_i,\quad r\to d_i+0:
  &&\sigma_\theta(r,\beta_i)
    \sim \frac{k_i}{\sqrt{2\pi(r-d_i)}},\\[3pt]
  &\theta=\beta_i,\quad r\to d_i-0:
  &&\frac{\partial u_\theta}{\partial r}(r,\beta_i)
    \sim -\frac{2(1-\nu_i^2)}{E_i}
    \frac{k_i}{\sqrt{2\pi(d_i-r)}}.
\end{aligned}
\end{equation}
Here, $k_i$ is the local Mode~I stress intensity factor at the process-zone
tip; $r\to d_i+0$ and $r\to d_i-0$ denote approach from ahead of the
zone and from within it, respectively. The boundary-value problem is closed by
imposing the standard stress-regularity condition for this
model~\citep{KaminskyDudykPolishchuk2024},
\begin{equation*}
  k_i=0.
\end{equation*}
This condition removes the inverse-square-root singular term and is
used below to determine the equilibrium zone length $d_i$. It does not
require the bounded part of the stress at the zone tip to vanish.


\section{Determination of the equilibrium process-zone parameters}
\label{sec:derivation}

\subsection{Mellin transform and functional equation}
\label{subsec:mellin}

To solve the local problem, we introduce the dimensionless radial coordinate
\begin{equation*}
  \rho=\frac{r}{d_i}
\end{equation*}
and apply the Mellin integral transform to the boundary-value problem
\eqref{eq:interface-boundary-conditions}--\eqref{eq:matched-far-field}.
Let $p$ be the complex Mellin variable. We define the transforms of the
correction normal stress outside the zone and the normal displacement
gradient within the zone by
\begin{equation*}
\begin{aligned}
  \Phi_i^+(p)
    &=\int_1^\infty
      \widehat{\sigma}_i(\rho d_i)\rho^p\,\dd\rho,\\
  \Phi_i^-(p)
    &=\frac{E_i}{2(1-\nu_i^2)}
      \int_0^1
      \left.
      \frac{\partial u_\theta}{\partial r}
      \right|_{\substack{r=\rho d_i\\\theta=\beta_i}}
      \rho^p\,\dd\rho.
\end{aligned}
\end{equation*}
Because $\widehat{\sigma}_i(r)=o(r^{-1})$ as $r\to\infty$, the integral
for $\Phi_i^+$ defines an analytic function for
$\operatorname{Re}p<0$. In the narrower region
$\operatorname{Re}p<-1-\lambda$, where the raw total-stress transform
\begin{equation*}
  \widetilde{\Phi}_i^+(p)
  =\int_1^\infty
    \sigma_\theta(\rho d_i,\beta_i)\rho^p\,\dd\rho
\end{equation*}
also converges, the two functions satisfy
\begin{equation*}
  \Phi_i^+(p)
  =\widetilde{\Phi}_i^+(p)
   +\frac{CQ_i d_i^\lambda}{p+\lambda+1},
\end{equation*}
thereby removing the pole generated by the matched outer field. Near the
corner, the displacement gradient scales as $r^{\lambda_i}$, so
$\Phi_i^-$ is analytic for $\operatorname{Re}p>-1-\lambda_i$. The superscripts on $\Phi_i^+$ and
$\Phi_i^-$ identify functions assigned to the left and right half-planes,
respectively; they do not indicate the sign of a physical quantity.

On the active segment $0<\rho<1$, the correction traction is
$\sigma_i-CQ_i d_i^\lambda\rho^\lambda$, and hence
\begin{equation*}
  \int_0^1
  \left(\sigma_i-CQ_i d_i^\lambda\rho^\lambda\right)
  \rho^p\,\dd\rho
  =
  \frac{\sigma_i}{p+1}
  -\frac{CQ_i d_i^\lambda}{p+\lambda+1}.
\end{equation*}
Eliminating the transforms of the interior field yields the functional
equation
\begin{equation}
\label{eq:wiener-hopf-functional}
  \Phi_i^+(p)
  +\frac{\sigma_i}{p+1}
  -\frac{CQ_i d_i^\lambda}{p+\lambda+1}
  =-\tan(\pi p)G_i(p)\Phi_i^-(p),
\end{equation}
which is initially valid in the common convergence strip
\begin{equation*}
  \max\{-1-\lambda,-1-\lambda_i\}
  <\operatorname{Re}p<0.
\end{equation*}
Here, $\operatorname{Re}p$ denotes the real part of $p$. In this strip,
both sides represent the same analytic relation. The functional equation then provides the analytic (or, in the presence of
the explicit rational terms, meromorphic) continuation of the combinations
required for factorization to the contour $\operatorname{Re}p=0$. This
continuation does not require the integral defining $\Phi_i^+$ itself to
converge on the imaginary axis. The
kernel $G_i(p)$ is
\begin{equation*}
  G_i(p)=\frac{D_i(p)\cos(\pi p)}{D(p)\sin(\pi p)}.
\end{equation*}
Here, $D(p)$ is the characteristic determinant for the corner with the
prescribed stationary V-shaped interfacial shear crack, given in
Appendix~\ref{app:asymptotics}.
For plane strain, define the dimensionless elastic ratio $e$ and the
material parameters $\kappa_i$ by
\begin{equation*}
  e=\frac{1+\nu_2}{1+\nu_1}\frac{E_1}{E_2},
  \qquad
  \kappa_i=3-4\nu_i.
\end{equation*}
The characteristic determinants $D_1(p)$ and $D_2(p)$ for configurations
with a process zone in materials~1 and~2, respectively, are
\begin{equation*}
\begin{aligned}
  D_1(p)={}&2e(1+\kappa_2)
    [\cos(2p\alpha)-\cos(2\alpha)]
    \{\sin^2[p(\pi-\alpha)]-p^2\sin^2\alpha\}\\
  &+(1+\kappa_1)
    [\sin(2p\alpha)+p\sin(2\alpha)]
    [\sin(2p(\pi-\alpha))-p\sin(2\alpha)],\\[2pt]
  D_2(p)={}&2(1+\kappa_1)
    [\cos(2p(\pi-\alpha))-\cos(2\alpha)]
    [\sin^2(p\alpha)-p^2\sin^2\alpha]\\
  &+e(1+\kappa_2)
    [\sin(2p(\pi-\alpha))-p\sin(2\alpha)]
    [\sin(2p\alpha)+p\sin(2\alpha)].
\end{aligned}
\end{equation*}
In particular, the corner singularity exponent $\lambda_i$ in the
presence of a zone in material~$i$ is determined by
\begin{equation*}
  D_i(-1-\lambda_i)=0,
\end{equation*}
where $\lambda_i$ is the smallest real root in $(-1,0)$ continuously
connected to the corresponding physical branch. As before,
$\alpha=\pi/2$ is understood only as a degenerate limit.

\subsection{Wiener--Hopf factorization}
\label{subsec:wiener-hopf}

We solve Eq.~\eqref{eq:wiener-hopf-functional} by the Wiener--Hopf
method~\citep{noble1958wienerhopf}, following the approach used for a related
boundary-value problem~\citep{KaminskyDudykPolishchuk2024}.
The factors $G_i^\pm(p)$ are defined by the Cauchy integral
\begin{equation*}
  \exp\left\{
    \frac{1}{2\pi\ii}
    \int_{-\ii\infty}^{\ii\infty}
      \frac{\log G_i(z)}{z-p}\,\dd z
  \right\}
  =
  \begin{cases}
    G_i^+(p), & \operatorname{Re}p<0,\\
    G_i^-(p), & \operatorname{Re}p>0,
  \end{cases}
\end{equation*}
where $z$ is the complex integration variable and $\ii^2=-1$.
The contour is directed upward along the imaginary axis. We also introduce
the auxiliary gamma-function factors
\begin{equation*}
  K^\pm(p)=
  \frac{\Gamma(1\mp p)}{\Gamma(\tfrac12\mp p)}.
\end{equation*}
Here, $\Gamma$ is the Euler gamma function.
With the normalization $G_i^\pm(p)\to1$ at infinity, the solution of
the functional equation is
\begin{equation}
\label{eq:wiener-hopf-solution}
\begin{aligned}
  \Phi_i^+(p)={}&
  -\frac{pG_i^+(p)}{K^+(p)}
  \Bigg\{
    \frac{\sigma_i}{p+1}
    \left[
      \frac{K^+(p)}{pG_i^+(p)}
      +\frac{K^+(-1)}{G_i^+(-1)}
    \right]\\
  &\hspace{37mm}
    -\frac{CQ_i d_i^\lambda}{p+\lambda+1}
    \left[
      \frac{K^+(p)}{pG_i^+(p)}
      +\frac{K^+(-\lambda-1)}
      {(\lambda+1)G_i^+(-\lambda-1)}
    \right]
  \Bigg\},
  \qquad \operatorname{Re}p<0,\\[3pt]
  \Phi_i^-(p)={}&
  K^-(p)G_i^-(p)
  \Bigg\{
    \frac{\sigma_iK^+(-1)}{(p+1)G_i^+(-1)}
    -\frac{CQ_i d_i^\lambda K^+(-\lambda-1)}
    {(\lambda+1)(p+\lambda+1)G_i^+(-\lambda-1)}
  \Bigg\},
  \qquad \operatorname{Re}p>0.
\end{aligned}
\end{equation}
For the kernels considered here, numerical checks show that
$G_i(\ii t)$ is real, positive, and even in the real variable $t$ and
tends to unity as $|t|\to\infty$.
Thus, for real $p<0$, the plus factor can be evaluated numerically
using the real integral
\begin{equation*}
  \log G_i^+(p)
  =-\frac{p}{\pi}\int_0^\infty
    \frac{\log G_i(\ii t)}{t^2+p^2}\,\dd t,
  \qquad p<0.
\end{equation*}
In the numerical calculations, we verify that the kernel is real and
positive on the factorization contour and approaches unity for large
$|t|$.

\subsection{Process-zone length}
\label{subsec:zone-length}

As $p\to+\infty$, the second relation in
Eq.~\eqref{eq:wiener-hopf-solution} gives
\begin{equation*}
  \Phi_i^-(p)\sim\frac{1}{\sqrt{p}}
  \left[
    \frac{\sigma_iK^+(-1)}{G_i^+(-1)}
    -\frac{CQ_i d_i^\lambda K^+(-\lambda-1)}
    {(\lambda+1)G_i^+(-\lambda-1)}
  \right].
\end{equation*}
Applying an Abelian-type theorem~\citep{noble1958wienerhopf} to the
tip asymptotics~\eqref{eq:zone-tip-asymptotics} also gives
\begin{equation*}
  \Phi_i^-(p)\sim-\frac{k_i}{\sqrt{2pd_i}}.
\end{equation*}
Comparing these expressions yields the local stress intensity factor
\begin{equation*}
  k_i=-\sqrt{2d_i}
  \left[
    \frac{\sigma_iK^+(-1)}{G_i^+(-1)}
    -\frac{CQ_i d_i^\lambda K^+(-\lambda-1)}
    {(\lambda+1)G_i^+(-\lambda-1)}
  \right].
\end{equation*}
Imposing $k_i=0$, which removes the inverse-square-root singular term
from the assumed tip asymptotics, gives the zone length:
\begin{equation}
\label{eq:process-zone-length}
  d_i=
  \left[
    \frac{CQ_iK^+(-\lambda-1)G_i^+(-1)}
    {\sigma_i(1+\lambda)K^+(-1)G_i^+(-\lambda-1)}
  \right]^{-1/\lambda}.
\end{equation}
Since $-1<\lambda<0$, a real positive $d_i$ requires $CQ_i>0$,
consistent with the physical admissibility condition in
Section~\ref{subsec:boundary-conditions}. For fixed geometry and fixed
V-shaped crack-branch length $l$, Eq.~\eqref{eq:process-zone-length}
gives the power law
$d_i\propto |C|^{-1/\lambda}$. Thus, varying the magnitude of the
external loading $|C|$ generates a quasi-static family of equilibrium
process-zone states $d_i(C)$ while the length of each V-shaped crack
branch remains unchanged. In the idealized singular formulation, this
relation does not itself define a finite threshold for formation of the
zone: every nonzero physically admissible $C$ gives $d_i>0$, with
$d_i\to0$ as $|C|\to0$. A finite material threshold enters at the next
stage, when branching initiation is identified from a critical
characteristic of the equilibrium zone.

At fixed geometry, V-shaped crack-branch length $l$, elastic properties,
and applied loading, the same equation gives the dependence on the normal cohesive
traction:
\begin{equation*}
  d_i\propto \sigma_i^{1/\lambda}
  =\sigma_i^{-1/|\lambda|}.
\end{equation*}
Because $-1<\lambda<0$, decreasing $\sigma_i$ increases the
process-zone length.

Using the relation
\begin{equation*}
  Q_i=g_1q_i l^{\lambda_0-\lambda},
\end{equation*}
where $g_1$ is the intact-corner angular coefficient and $q_i$ is a
dimensionless amplitude coefficient, both given in
Appendix~\ref{app:asymptotics}, we obtain from
Eq.~\eqref{eq:process-zone-length}
\begin{equation}
\label{eq:normalized-process-zone-length}
  \frac{d_i}{l}=
  \left[
    \frac{Cl^{\lambda_0}}{\sigma_i}
    \frac{g_1q_iK^+(-1-\lambda)G_i^+(-1)}
    {(1+\lambda)K^+(-1)G_i^+(-1-\lambda)}
  \right]^{-1/\lambda}.
\end{equation}
This relation also provides a check on the small-scale assumption
$d_i/l\ll1$ used in the local formulation.

\subsection{Opening and energy parameter}
\label{subsec:zone-parameters}

Let the process-zone opening at the corner be
$\delta_i=\jump{u_\theta(0,\beta_i)}$.
The definition of $\Phi_i^-(p)$ gives
\begin{equation*}
  \delta_i
  =-2
    \left.
      \left(\frac{\partial u_\theta}{\partial r}\right)^*
    \right|_{p=0,\,\theta=\beta_i}
  =-\frac{4(1-\nu_i^2)d_i}{E_i}\Phi_i^-(0).
\end{equation*}
Here, the asterisk denotes the dimensional transform moment:
for $f(r)=\partial u_\theta(r,\beta_i)/\partial r$,
$f^*(p)=\int_0^{d_i}f(r)(r/d_i)^p\,\dd r$.
The factor of two accounts for the symmetric displacement of the two
zone faces. Substitution of Eq.~\eqref{eq:wiener-hopf-solution},
together with Eq.~\eqref{eq:process-zone-length}, gives
\begin{equation}
\label{eq:process-zone-opening}
  \delta_i=
  -\frac{4(1-\nu_i^2)\sigma_i\lambda K^+(-1)}
  {E_i\sqrt{\pi G_i(0)}(1+\lambda)G_i^+(-1)}d_i.
\end{equation}
For fixed material properties and geometry, the coefficient multiplying
$d_i$ in Eq.~\eqref{eq:process-zone-opening} is constant. The opening
therefore scales with the zone length as the applied loading varies:
$\delta_i\propto d_i$. If the applied loading is fixed and $\sigma_i$
is varied, the preceding power law instead gives
\begin{equation*}
  \delta_i\propto \sigma_i d_i
  \propto \sigma_i^{1+1/\lambda}
  =\sigma_i^{1-1/|\lambda|}.
\end{equation*}
The exponent $1-1/|\lambda|$ is negative. A decrease in the resistance
to separation therefore increases both the zone length and its opening.

Within the adopted model, the work $W_i$ done against the resistance
to separation, per unit out-of-plane thickness, is defined
by~\citep{Cherepanov1986}
\begin{equation*}
  W_i=\sigma_i\int_0^{d_i}
      \jump{u_\theta(r,\beta_i)}\,\dd r
  =-\frac{4(1-\nu_i^2)\sigma_i}{E_i}d_i^2\Phi_i^-(1).
\end{equation*}
Using Eq.~\eqref{eq:wiener-hopf-solution}, we obtain
\begin{equation*}
  W_i=-\frac{8\lambda(1-\nu_i^2)\sigma_i^2}
  {\pi(\lambda+2)[G_i^+(-1)]^2E_i}d_i^2.
\end{equation*}
Differentiating with respect to the zone length along the family of
equilibrium states, with the geometry and material properties fixed,
defines the zone-related energy parameter $J_i$:
\begin{equation}
\label{eq:process-zone-energy-parameter}
  J_i=\frac{\dd W_i}{\dd d_i}
  =-\frac{16\lambda(1-\nu_i^2)\sigma_i^2}
  {\pi(\lambda+2)[G_i^+(-1)]^2E_i}d_i.
\end{equation}
For the numerical analysis, introduce the dimensionless loading,
opening, and energy parameters $\sigmap$, $\delta_i'$, and $J_i'$,
respectively:
\begin{equation*}
  \sigmap=\frac{Cl^{\lambda_0}}{\sigma_i},
  \qquad
  \delta_i'=\frac{E_i\delta_i}
  {4(1-\nu_i^2)\sigma_i l},
  \qquad
  J_i'=\frac{E_iJ_i}
  {4(1-\nu_i^2)\sigma_i^2l}.
\end{equation*}
The primes here denote normalization, not differentiation.
Equations~\eqref{eq:normalized-process-zone-length},
\eqref{eq:process-zone-opening}, and
\eqref{eq:process-zone-energy-parameter} then take the form
\begin{equation*}
\begin{aligned}
  \frac{d_i}{l}
    &=\left[
      \sigmap\,
      \frac{g_1q_iK^+(-1-\lambda)G_i^+(-1)}
      {(1+\lambda)K^+(-1)G_i^+(-1-\lambda)}
    \right]^{-1/\lambda},\\[2pt]
  \delta_i'
    &=-\frac{\lambda K^+(-1)}
      {\sqrt{\pi G_i(0)}(1+\lambda)G_i^+(-1)}
      \frac{d_i}{l},\\[2pt]
  J_i'
    &=-\frac{4\lambda}
      {\pi(\lambda+2)[G_i^+(-1)]^2}
      \frac{d_i}{l}.
\end{aligned}
\end{equation*}
All three relations use the corner exponent $\lambda$ for the prescribed
stationary V-shaped interfacial shear crack, i.e., the root of
$D(-1-\lambda)=0$.
The exponents $\lambda_1$ and $\lambda_2$, which describe the corner
singularity in the presence of a zone, do not replace $\lambda$
in the formulas for the zone parameters.

For fixed geometry and material properties, both $\delta_i$ and $J_i$
depend linearly on the same equilibrium zone length $d_i$. Eliminating
$d_i$ between Eqs.~\eqref{eq:process-zone-opening} and
\eqref{eq:process-zone-energy-parameter} gives
\begin{equation}
\label{eq:energy-opening-relation}
  J_i=\chi_i\sigma_i\delta_i,
  \qquad
  \chi_i=
  \frac{4(1+\lambda)\sqrt{\pi G_i(0)}}
  {\pi(\lambda+2)G_i^+(-1)K^+(-1)}.
\end{equation}
Here, $\chi_i$ is a dimensionless proportionality factor determined by
the local geometry and elastic contrast.

Within the adopted local model, $J_i$ is interpreted as a zone-related
energy parameter rather than as an independently established global contour
integral. In analogy with process-zone criteria used for interface-crack
kinking~\citep{KaminskyEtAl2023}, branching initiation may be expressed
through the energy criterion
\begin{equation*}
  J_i=J_{i\mathrm{c}},
\end{equation*}
where $J_{i\mathrm{c}}$ is an independently specified critical
characteristic within this criterion. Alternatively, a deformation criterion
\begin{equation*}
  \delta_i=\delta_{i\mathrm{c}}
\end{equation*}
may be used, with $\delta_{i\mathrm{c}}$ specified independently as the
critical opening. For fixed geometry and material properties,
Eq.~\eqref{eq:energy-opening-relation} shows that the two criteria
describe the same one-parameter family of local states and select the
same critical state only when the independently specified critical
parameters are calibrated consistently:
\begin{equation*}
  J_{i\mathrm{c}}
  =\chi_i\sigma_i\delta_{i\mathrm{c}}.
\end{equation*}


\section{Numerical results}
\label{sec:numerical-results}

The numerical analysis is performed for the normalized parameters
$\,d_i/l$, $\delta_i'$, and $J_i'$ defined in
Section~\ref{sec:derivation}. Only combinations of material, angle, and
loading sign for which
\begin{equation*}
  Cg_2<0,\qquad CQ_i>0
\end{equation*}
hold simultaneously are considered physically admissible. The first
inequality ensures compressive contact between the faces of the V-shaped
interfacial shear crack, whereas the second ensures tensile normal stress
along the bisector of the angular sector of the material in which the
process zone forms.

For the baseline material pair
\begin{equation*}
  E_1/E_2=0.5,\qquad \nu_1=\nu_2=0.3,
\end{equation*}
the sign changes of the coefficient $g_2$ occur at
\begin{equation*}
  \alpha_1=12.92^\circ,\qquad
  \alpha_2=107.09^\circ.
\end{equation*}
These values are obtained by directly solving
$D_0(-1-\lambda_0)=0$ together with $g_2=0$. For $C>0$, a process zone
is physically admissible in material~2 for $0<\alpha<\alpha_1$ and in
material~1 for $\alpha_2<\alpha<180^\circ$. For $C<0$, the corresponding
intervals are $\alpha_1<\alpha<90^\circ$ for material~1 and
$90^\circ<\alpha<\alpha_2$ for material~2. The value
$\alpha=90^\circ$ corresponds to a straight interface and is treated
below as a degenerate analytical limit rather than as an ordinary
interior point of a numerical branch.

To assess the effect of elastic contrast, the same two transition angles
were calculated for a range of $E_1/E_2<1$ with
$\nu_1=\nu_2=0.3$. The results are given in
Table~\ref{tab:g2-transition-angles}. As $E_1/E_2$ increases, the first
transition angle $\alpha_1$ increases monotonically, whereas $\alpha_2$
decreases monotonically. Thus, elastic contrast systematically shifts
the boundaries of the angular intervals of physical admissibility.

\begin{table}[htbp]
  \centering
  \caption{Angles at which the coefficient $g_2$ changes sign for
  $\nu_1=\nu_2=0.3$.}
  \label{tab:g2-transition-angles}
  \scriptsize
  \setlength{\tabcolsep}{2.6pt}
  \begin{tabular}{@{}lrrrrrrrrrrr@{}}
    \toprule
    $E_1/E_2$ & 0.01 & 0.10 & 0.20 & 0.30 & 0.40 & 0.50 & 0.60 & 0.70 & 0.80 & 0.90 & 0.99 \\
    \midrule
    $\alpha_1$, $^\circ$ & 2.58 & 7.03 & 9.25 & 10.77 & 11.95 & 12.92 & 13.74 & 14.45 & 15.08 & 15.65 & 16.11 \\
    $\alpha_2$, $^\circ$ & 108.02 & 107.90 & 107.72 & 107.52 & 107.30 & 107.09 & 106.88 & 106.69 & 106.50 & 106.32 & 106.17 \\
    \bottomrule
  \end{tabular}
\end{table}


Cases with $E_1/E_2>1$ are not independent. They are obtained by
interchanging the materials according to
\begin{equation*}
  (E_1,\nu_1)\leftrightarrow(E_2,\nu_2),
  \qquad
  \alpha\leftrightarrow\pi-\alpha,
  \qquad
  1\leftrightarrow2.
\end{equation*}
For the corresponding local quantities, this mapping also gives
$d_1\leftrightarrow d_2$ and
$\lambda_1\leftrightarrow\lambda_2$; for the transition half-angles,
$\alpha_1\leftrightarrow\pi-\alpha_2$. For $\nu_1=\nu_2$, this
includes the replacement
$E_1/E_2\leftrightarrow E_2/E_1$, so Table~\ref{tab:g2-transition-angles}
need only list $E_1/E_2<1$.

\subsection{Dependence on normalized loading}
\label{subsec:figure2}

Figure~\ref{fig:load-dependence} shows the process-zone parameters as
functions of
\begin{equation*}
  \sigma'=\frac{Cl^{\lambda_0}}{\sigma_i}
\end{equation*}
for four reference configurations:
$\alpha=10^\circ$ ($i=2$, $C>0$),
$45^\circ$ ($i=1$, $C<0$),
$105^\circ$ ($i=2$, $C<0$), and
$135^\circ$ ($i=1$, $C>0$). For each configuration, only the physically
admissible half-axis of the loading parameter is shown.

\resultfigure{Fig2.pdf}{%
Dependence of the normalized process-zone parameters on the
dimensionless loading parameter $\sigma'$ for four reference angles and
the physically admissible materials.}{fig:load-dependence}

In all four cases, $d_i/l$, $\delta_i'$, and $J_i'$ increase
monotonically with $|\sigma'|$. This common trend follows directly from
the analytical relations
$\delta_i'\propto d_i/l$ and $J_i'\propto d_i/l$ for fixed geometry.
Among the four reference configurations shown, the largest values of all
three parameters at $|\sigma'|=0.5$ occur for
$\alpha=45^\circ$, $C<0$; specifically,
\begin{equation*}
  \frac{d_1}{l}=0.1474,\qquad
  \delta_1'=0.1458,\qquad
  J_1'=0.0836.
\end{equation*}
This trend also bears on the small-scale assumption: increasing the load
raises the local deformation and energy parameters and moves the model
closer to the limit of its asymptotic applicability.

These loading curves also provide a direct procedure for determining
the critical load for branching initiation once a critical material
parameter is specified independently. If the material is characterized by
a critical opening $\delta_{i\mathrm{c}}$ or a critical zone-related
energy parameter $J_{i\mathrm{c}}$, the intersection with the
corresponding curve determines the normalized critical load $\sigma'_c$
and hence the amplitude of the outer field,
\begin{equation*}
  C_c=\sigma'_c\sigma_i l^{-\lambda_0}.
\end{equation*}
No numerical value of $\delta_{i\mathrm{c}}$ or $J_{i\mathrm{c}}$ is
prescribed in the present examples, so the plots do not identify a
particular initiation point. Equation~\eqref{eq:energy-opening-relation}
shows that, for fixed geometry and material properties, $\delta_i$ and
$J_i$ are linearly related measures of the same equilibrium zone length.
Accordingly, when
$J_{i\mathrm{c}}=\chi_i\sigma_i\delta_{i\mathrm{c}}$, the
deformation and energy criteria identify the same critical state.

\subsection{Dependence on the interface half-angle}
\label{subsec:figure3}

Figure~\ref{fig:angle-dependence} shows the dependence of the normalized
process-zone parameters on the half-angle $\alpha$ for the same baseline
material pair at $\sigma'=\pm0.5$. Only material/loading-sign
combinations satisfying the physical-admissibility conditions are
retained. Transitions between admissible branches occur at
$\alpha_1$, $90^\circ$, and $\alpha_2$.

\resultfigure{Fig3.pdf}{%
Dependence of the normalized process-zone parameters on the half-angle
$\alpha$ at $\sigma'=\pm0.5$. Only physically admissible combinations
of material and loading sign are shown; $\alpha=90^\circ$ is an
analytical limit.}%
{fig:angle-dependence}

On the two broad admissible material-1 branches, the angular dependencies
have pronounced interior maxima, whereas on the short material-2 branches
the largest values occur near the admissibility boundaries. The maxima of $d_i/l$, $\delta_i'$, and $J_i'$ are close but do not occur
at exactly the same angle. For the material-1 branch with $C<0$, the
one-degree sweep gives
\begin{equation*}
\begin{aligned}
  \max(d_1/l)&=0.1779 &&\text{at }\alpha=35^\circ,\\
  \max\delta_1'&=0.2077 &&\text{at }\alpha=32^\circ,\\
  \max J_1'&=0.1139 &&\text{at }\alpha=33^\circ.
\end{aligned}
\end{equation*}
Their proximity is consistent with the fact that, at each fixed $\alpha$,
$\delta_i'$ and $J_i'$ are proportional to $d_i/l$, while the
proportionality coefficients themselves depend on angle. Thus,
Figure~\ref{fig:angle-dependence} quantifies how geometry changes the
deformation and energy measures attained at a prescribed load. Converting
these curves into an angular dependence of the critical applied load requires
specifying the corresponding critical material parameter and solving the
criterion separately for each $\alpha$.

The maximum of the material-1 branch with $C<0$ in
Figure~\ref{fig:angle-dependence} corresponds to $d_i/l$ of order
$0.18$. This value is no longer unambiguously small compared with unity,
whereas the local formulation is based on the assumption
$d_i/l\ll1$. We therefore use $d_i/l<0.1$ as the principal range for
quantitative interpretation, where the small-scale approximation is best
justified. Portions of the curves with $d_i/l>0.1$ are retained to show
the structure and trends of the analytical solution, but they should be
viewed as extrapolations of the local model rather than equally accurate
quantitative predictions.

Numerical evaluation of the plus factors for
Figure~\ref{fig:angle-dependence} requires particular care as
$\alpha\to0$ and $\alpha\to180^\circ$. A fixed factorization contour
converges slowly near these limits. The calculations therefore use
adaptive contour extension; a further extension of the tail changed the
computed plus factors by at most $8.69\times10^{-13}$. Accordingly,
Figure~\ref{fig:angle-dependence} is based on the converged adaptive
values.

\subsection{Singularity exponents and physical admissibility}
\label{subsec:figure4}

To compare the local stress-field structure, consider three mathematical
configurations of the same corner region: the intact bimaterial corner,
the corner containing the prescribed stationary interfacial shear crack,
and the configuration containing both that crack and a local process zone.
For the baseline material pair, the corresponding singularity orders are
described by the exponents $\lambda_0$, $\lambda$, and $\lambda_i$.
This comparison characterizes differences among the local field models
and does not prescribe a chronological sequence of fracture development.

\resultfigure{Fig4.pdf}{%
Singularity exponents and physically admissible segments for the
baseline material pair. Complete mathematical branches and physically
admissible portions are distinguished according to the conditions
$Cg_2<0$ and $CQ_i>0$.}%
{fig:singularity-exponents}

For the baseline material pair, within the physically admissible segments,
the configuration containing the interfacial shear crack has a substantially
stronger singularity than the intact corner:
$|\lambda|$ exceeds $|\lambda_0|$. Introducing the process zone partly weakens this singularity, and for the
corresponding material
\begin{equation*}
  |\lambda_0|<|\lambda_i|<|\lambda|.
\end{equation*}
Thus, the process zone has a local shielding effect: it reduces the
order of the singularity associated with the interfacial shear crack but
does not eliminate the stress concentration completely. In addition, on
all physically admissible open intervals of the baseline calculation,
\begin{equation*}
  -\frac12<\lambda_i<0,
\end{equation*}
so the residual singularity at $O$ is weaker than the classical
$r^{-1/2}$ crack singularity. Both results are established here for
the material pair considered and its physically admissible branches and
should not be extended to arbitrary elastic contrasts without additional
verification.

The stronger singularity of the configuration containing the interfacial
crack creates conditions favorable to local fracture in one of the
materials, but the existence and size of a process zone also depend on
the conditions $Cg_2<0$ and $CQ_i>0$, the resistance to separation
$\sigma_i$, and the level of external loading. The exponents
$\lambda_0$, $\lambda$, and $\lambda_i$ therefore characterize the
strengthening and shielding of the local field, whereas branching
initiation is associated with attainment of a critical $\delta_i$ or
$J_i$ along the corresponding family of equilibrium states.


\section{Discussion}
\label{sec:discussion}

\subsection{Quasi-static family of states and branching initiation}
\label{subsec:discussion-quasistatic}

In the present formulation, the V-shaped interfacial shear crack is a
prescribed stationary configuration: the length $l$ of each crack branch
remains fixed, whereas the equilibrium process-zone length $d_i$ changes
with the outer-loading parameter $C$. The solution therefore defines the
one-parameter quasi-static family
\begin{equation*}
  C\longrightarrow d_i(C)
  \longrightarrow
  \bigl\{\delta_i(C),W_i(C),J_i(C)\bigr\}.
\end{equation*}
Here, quasi-static refers to a family of equilibrium zone configurations
in which inertia is neglected.

This family can be used directly to determine the critical load for
branching initiation. If a critical opening $\delta_{i\mathrm{c}}$ or a
critical zone-related energy parameter $J_{i\mathrm{c}}$ is specified
independently for the material, the critical state is defined by
\begin{equation*}
  \delta_i(C_c)=\delta_{i\mathrm{c}}
  \qquad\text{or}\qquad
  J_i(C_c)=J_{i\mathrm{c}}.
\end{equation*}
Once the corresponding critical material characteristic is known, these
relations determine $C_c$. Subsequent propagation and the finite path of
the branched crack require a separate formulation.

\subsection{Competing fracture mechanisms}
\label{subsec:discussion-competing}

Two distinct local mechanisms are possible for the stationary V-shaped
crack. The first is associated with its corner $O$: the singular corner
field supports a tensile process zone in one of the materials, and
attainment of a critical $\delta_i$ or $J_i$ marks branching initiation.
The second mechanism is localized at the tips of the V-shaped crack branches.
For those tips, \citet{NazarenkoKipnis2022Equilibrium} determined the Mode~II stress
intensity factor $K_{\mathrm{II}}$ and formulated the interfacial
fracture condition
\begin{equation*}
  K_{\mathrm{II}}=K_{\mathrm{IIc}}.
\end{equation*}

Branching from the corner and interfacial propagation from the crack-branch tips
should therefore be regarded as competing local mechanisms of the same
stationary geometry, rather than as prescribed sequential stages. Their
respective critical loads are governed by different material
characteristics: Mode~II interfacial fracture resistance at the crack-branch tips
and the critical process-zone characteristic, $\delta_i$ or $J_i$, at
the corner.

The small-scale assumption $d_i\ll l$ permits the two local regions to be
treated separately at leading order. The field of the stationary V-shaped
crack serves as the outer field for the process-zone problem near $O$, and
the feedback of the small zone on the fields at the crack-branch tips is
neglected. A coupled boundary-value problem would be required for a
quantitative analysis of the competition when $d_i/l$ is no longer small.

\subsection{Energy and deformation criteria}
\label{subsec:discussion-criteria}

According to Eq.~\eqref{eq:energy-opening-relation}, for fixed geometry
and elastic properties,
\begin{equation*}
  J_i=\chi_i\sigma_i\delta_i.
\end{equation*}
Thus, $J_i$ and $\delta_i$ are not independent measures of the current
zone state: both are uniquely determined by the same equilibrium length
$d_i$. The deformation criterion
$\delta_i=\delta_{i\mathrm{c}}$ and the energy criterion
$J_i=J_{i\mathrm{c}}$ therefore select the same critical state when
\begin{equation*}
  J_{i\mathrm{c}}
  =\chi_i\sigma_i\delta_{i\mathrm{c}}.
\end{equation*}

This equivalence is specific to the present one-parameter model and fixed
geometry. It does not imply a universal identity between two independently
determined critical characteristics for arbitrary fracture problems, because
$\chi_i$ depends on the corner geometry and elastic contrast. In
applications, either critical characteristic may be used, provided it is
established independently and the two criteria are not treated as
independent constraints on the same model state.

\subsection{Residual corner field and local energy flux}
\label{subsec:discussion-flux}

In the configuration containing the process zone, the field at the corner
$O$ remains singular. If $H_i$ is the amplitude of the leading corner
eigenfield, then
\begin{equation*}
  \sigma\sim H_i r^{\lambda_i}.
\end{equation*}
The corresponding elastic strain-energy density is of order
\begin{equation*}
  w\sim\frac{H_i^2}{E_i}r^{2\lambda_i},
\end{equation*}
and the local configurational energy flux through a contour of
characteristic radius $R$ scales as
\begin{equation}
\label{eq:local-elastic-flux-scaling}
  J_i^{(\mathrm{el})}(R)
  \sim\frac{H_i^2}{E_i}R^{2\lambda_i+1}.
\end{equation}

For the physically admissible regimes identified in the numerical
analysis, $-1/2<\lambda_i<0$. Hence $2\lambda_i+1>0$ and
\begin{equation*}
  \lim_{R\to0}J_i^{(\mathrm{el})}(R)=0.
\end{equation*}
The residual stress singularity at $O$ therefore does not generate a
finite nonzero local elastic contribution to the zone energy parameter.
In this sense, the corner singularity in the configuration containing the
process zone is weaker than the classical $r^{-1/2}$ crack singularity.

This distinction also identifies the limit of the additive representation
used by \citet{KaminskyDudyk2026Criteria}. For the
existing crack considered there, the outer field has a singularity exponent
with real part $-1/2$, and a finite elastic energy flux can enter the overall energy
balance together with the process-zone contribution. For the present
corner problem, $\lambda_i>-1/2$ implies that
$J_i^{(\mathrm{el})}(R)\to0$, so no separate finite elastic term
$J_i^{(\mathrm{el})}$ is required in the branching-initiation criterion. The zone-related parameter
$J_i=\dd W_i/\dd d_i$ can therefore be used without such an additive
correction; this does not make $J_i$ an independently established global
contour integral.

\subsection{Range of applicability and generalization}
\label{subsec:discussion-scope}

The analysis is based on the scale separation $d_i\ll l\ll L$. We use
$d_i/l<0.1$ as the principal range for quantitative application. Portions
of the numerical curves for which $d_i/l$ exceeds $0.1$ and reaches
approximately $0.18$ near the maxima remain useful for showing trends, but
they represent extrapolations of the local approximation. In those
regimes, higher-order terms of the outer asymptotic field and interaction
between the local regions near the corner and the crack-branch tips become
increasingly important.

The local character of the solution also facilitates extension to a
broader class of problems. For a developed symmetric V-shaped interfacial
crack with finite branches, the factor $CQ_i$ that supplies the outer-field
amplitude in the local problem can be replaced by the amplitude, or
generalized stress intensity factor, of the corresponding corner singular
field. Such an extension remains consistent provided that the process zone
is small and the same local eigenproblem applies. The model can therefore
be used to analyze branching initiation at a kink in an interface. A
related application is splitting of a solid by a sharp wedge when its
local boundary-value problem reduces to an analogous corner configuration.


\section{Conclusions}
\label{sec:conclusions}

\begin{enumerate}
\item
An analytical model has been developed for branching initiation at the
corner of a V-shaped interfacial shear crack. The singular field of the
prescribed stationary crack is used as the outer field for a small tensile
process zone oriented along the bisector of one of the materials.
Closed-form relations are obtained for the equilibrium zone length $d_i$,
opening $\delta_i$, work $W_i$ done against the resistance to separation,
and the zone-related energy parameter $J_i=\dd W_i/\dd d_i$.

\item
Physical admissibility of the local process zone requires both compressive
contact of the V-shaped interfacial crack faces, $Cg_2<0$, and a tensile
normal field in the material containing the zone, $CQ_i>0$. For the
parameters examined, these conditions can be satisfied simultaneously in
only one of the two materials for a given angle and loading sign. The
interface angle and elastic contrast shift the admissible ranges and the
equilibrium zone characteristics.

\item
The presence of the process zone weakens the corner singularity of the
configuration but does not eliminate it. On the physically admissible intervals for the baseline
material pair, $-1/2<\lambda_i<0$. The residual field is therefore weaker
than the classical $r^{-1/2}$ crack singularity, and its local
configurational elastic energy flux vanishes as the contour contracts to
the corner. Consequently, no separate finite elastic additive term is required for the
zone-related energy parameter $J_i$ in the local branching-initiation
criterion for this corner problem.

\item
For fixed geometry, the zone opening and the zone-related energy parameter
are uniquely related through the same equilibrium zone length:
$J_i=\chi_i\sigma_i\delta_i$. The deformation criterion
$\delta_i=\delta_{i\mathrm{c}}$ and the energy criterion
$J_i=J_{i\mathrm{c}}$ therefore identify the same critical state only
when the independently specified critical parameters are calibrated
consistently,
$J_{i\mathrm{c}}=\chi_i\sigma_i\delta_{i\mathrm{c}}$. Once either
critical material parameter is supplied independently, the model determines
the corresponding critical external load for branching initiation.

\item
Quantitative use of the model relies on the small-scale assumption
$d_i\ll l$; $d_i/l<0.1$ is adopted as the principal applicability range,
whereas portions of the curves reaching approximately $d_i/l=0.18$ should
be interpreted as extrapolations of the local approximation. Branching
from the corner and interfacial propagation from the tips of the V-shaped crack
branches are competing local mechanisms rather than prescribed
sequential stages. Because the formulation is local, it can be extended
to developed V-shaped interfacial cracks by replacing the outer-field
amplitude with the corresponding corner-field amplitude. A related
application is splitting of a solid by a sharp wedge.
\end{enumerate}


\begin{appendices}
\section{Stress and displacement fields near the corner}
\label{app:asymptotics}
\providecommand{\SingularityExponentFigure}{Figure~\ref{fig:singularity-exponents}}

\setcounter{equation}{0}
\renewcommand{\theequation}{A.\arabic{equation}}

This appendix collects the characteristic relations and asymptotic-field
coefficients used in the main text. For plane strain, the dimensionless
elastic ratio $e$ and material parameters $\kappa_i$ are
\begin{equation*}
  e=\frac{1+\nu_2}{1+\nu_1}\frac{E_1}{E_2},
  \qquad
  \kappa_i=3-4\nu_i,\qquad i=1,2.
\end{equation*}
For each characteristic exponent, the complex Mellin variable $p$ and
the generic singularity exponent $\lambda_*$ are related by
\begin{equation*}
  p=-1-\lambda_*,
  \qquad -1<\lambda_*<0.
\end{equation*}
The physical exponent is taken to be the smallest real root in
$(-1,0)$ that is continuously connected to the corresponding solution
branch. The case $\alpha=\pi/2$ is the degenerate straight-interface
limit and is not treated as an ordinary interior point of a
characteristic branch.

\subsection{Intact corner}

For the configuration without an interfacial shear crack, the singular
part of the interfacial stress field is~\citep{NazarenkoKipnis2022Stress}
\begin{equation}
\label{eq:app-intact-field}
\begin{aligned}
  \tau_{r\theta}(r,\alpha)
    &=Cg_1r^{\lambda_0}+o\!\left(r^{\lambda_0}\right),\\
  \sigma_\theta(r,\alpha)
    &=Cg_2r^{\lambda_0}+o\!\left(r^{\lambda_0}\right),
    \qquad r\to0,
\end{aligned}
\end{equation}
where $C$ is the local singular-field amplitude, $\lambda_0$ is
the intact-corner singularity exponent, and the dimensionless angular
coefficients $g_1$ and $g_2$ are
\begin{equation*}
\begin{aligned}
  g_1={}&(1-e)\sin[\lambda_0(\pi-\alpha)]
  \bigl[(\lambda_0+1)\sin2\alpha
       +\sin\{2(\lambda_0+1)\alpha\}\bigr]\\
  &+e(1+\kappa_2)\sin(\lambda_0\pi)
      \sin[(\lambda_0+2)\alpha],\\[2pt]
  g_2={}&(1-e)\cos[\lambda_0(\pi-\alpha)]
  \bigl[(\lambda_0+1)\sin2\alpha
       +\sin\{2(\lambda_0+1)\alpha\}\bigr]\\
  &+e(1+\kappa_2)\sin(\lambda_0\pi)
      \cos[(\lambda_0+2)\alpha].
\end{aligned}
\end{equation*}
The exponent $\lambda_0$ is determined by the characteristic equation
\begin{equation}
\label{eq:app-D0-root}
  D_0(-1-\lambda_0)=0,
  \qquad -1<\lambda_0<0,
\end{equation}
where the intact-corner determinant $D_0(p)$ is
\begin{equation*}
  D_0(p)=b_0(p)+e\,b_1(p)+e^2b_2(p),
\end{equation*}
with auxiliary functions $b_0(p)$, $b_1(p)$, and $b_2(p)$ given by
\begin{equation*}
\begin{aligned}
  b_0(p)={}&
  [\sin(2p\alpha)+p\sin(2\alpha)]
  [\kappa_1\sin\{2p(\pi-\alpha)\}+p\sin(2\alpha)],\\[2pt]
  b_1(p)={}&(1+\kappa_1)(1+\kappa_2)\sin^2(p\pi)\\
  &-[\sin(2p\alpha)+p\sin(2\alpha)]
    [\kappa_1\sin\{2p(\pi-\alpha)\}+p\sin(2\alpha)]\\
  &-[\sin\{2p(\pi-\alpha)\}-p\sin(2\alpha)]
    [\kappa_2\sin(2p\alpha)-p\sin(2\alpha)],\\[2pt]
  b_2(p)={}&
  [\sin\{2p(\pi-\alpha)\}-p\sin(2\alpha)]
  [\kappa_2\sin(2p\alpha)-p\sin(2\alpha)].
\end{aligned}
\end{equation*}
\subsection{Corner with a stationary V-shaped interfacial shear crack}

For the prescribed stationary V-shaped interfacial shear crack, the two
symmetric branches have length $l$ each. At distances $r\ll l$, the normal stress along the bisector
of the sector occupied by material~$i$ has the asymptotic
form~\citep{NazarenkoKipnis2022Stress}
\begin{equation}
\label{eq:app-cracked-field}
  \sigma_\theta(r,\beta_i)
  =CQ_i r^\lambda+o\!\left(r^\lambda\right),
  \qquad r\to0,
  \qquad
  \beta_1=\pi,\quad \beta_2=0.
\end{equation}
Here, $\lambda$ is the corner singularity exponent for the cracked
configuration and $\beta_i$ is the corresponding bisector direction.
The amplitude coefficient $Q_i$ is expressed as
\begin{equation*}
  Q_i=g_1q_i l^{\lambda_0-\lambda},
\end{equation*}
where the dimensionless coefficient $q_i$ and the derivative $D'(p)$
of the cracked-corner determinant are
\begin{equation*}
  q_i=
  \frac{2S_i(-1-\lambda)K^+(-1-\lambda)G^+(-1-\lambda)}
  {(\lambda-\lambda_0)D'(-1-\lambda)
   K^+(-1-\lambda_0)G^+(-1-\lambda_0)},
  \qquad
  D'(p)=\frac{\dd D(p)}{\dd p}.
\end{equation*}
The functions $S_i(p)$ are auxiliary numerator functions defined below.
The factors $K^+(p)$ and $G^+(p)$ are specified later in this subsection.
The base plus factor $G^+(p)$ must not be confused with $G_i^+(p)$, which
factorizes the process-zone kernel $G_i(p)$ in
Section~\ref{sec:derivation}. Explicitly,
\begin{equation*}
\begin{aligned}
  S_1(p)={}&2e(1+\kappa_2)s_{11}(p)\sin[(p+1)\alpha]\\
  &+s_{12}(p)
  \Bigl[(1+\kappa_1)\sin\{p(\pi-\alpha)-\alpha\}
        +2(1-e)s_{13}(p)\Bigr],\\[2pt]
  S_2(p)={}&s_{21}(p)
  \Bigl[e(1+\kappa_2)\sin[(p+1)\alpha]
        +2(1-e)s_{22}(p)\Bigr]\\
  &-2(1+\kappa_1)\sin\{p(\pi-\alpha)-\alpha\}\,s_{23}(p),
\end{aligned}
\end{equation*}
where the auxiliary functions $s_{mn}(p)$, with $m=1,2$ and
$n=1,2,3$, are
\begin{equation*}
\begin{aligned}
  s_{11}(p)&=p\cos(p\pi-\alpha)\sin\alpha
             -\sin[p(\pi-\alpha)]\cos(p\alpha),\\
  s_{12}(p)&=p\sin2\alpha+\sin(2p\alpha),\\
  s_{13}(p)&=p\cos[p(\pi-\alpha)]\sin\alpha
             -\sin[p(\pi-\alpha)]\cos\alpha,\\
  s_{21}(p)&=\sin[2p(\pi-\alpha)]-p\sin2\alpha,\\
  s_{22}(p)&=p\cos(p\alpha)\sin\alpha
             +\sin(p\alpha)\cos\alpha,\\
  s_{23}(p)&=p\cos(p\pi+\alpha)\sin\alpha
             +\cos[p(\pi-\alpha)]\sin(p\alpha).
\end{aligned}
\end{equation*}
The corner singularity exponent $\lambda$ for the configuration
containing the prescribed stationary V-shaped interfacial crack satisfies
\begin{equation}
\label{eq:app-D-root}
  D(-1-\lambda)=0,
  \qquad -1<\lambda<0,
\end{equation}
where the characteristic determinant $D(p)$ is
\begin{equation*}
  D(p)=a_0(p)+e\,a_1(p),
\end{equation*}
with auxiliary functions $a_0(p)$ and $a_1(p)$ given by
\begin{equation*}
\begin{aligned}
  a_0(p)={}&(1+\kappa_1)
  [\cos\{2p(\pi-\alpha)\}-\cos2\alpha]
  [\sin(2p\alpha)+p\sin2\alpha],\\
  a_1(p)={}&(1+\kappa_2)
  [\cos(2p\alpha)-\cos2\alpha]
  [\sin\{2p(\pi-\alpha)\}-p\sin2\alpha].
\end{aligned}
\end{equation*}
The coefficients $q_i$ use the base kernel
\begin{equation*}
  G(p)=
  \frac{(e+\kappa_1)(1+e\kappa_2)\sin(\pi p)D(p)}
  {(e+\kappa_1+1+e\kappa_2)\cos(\pi p)D_0(p)}
\end{equation*}
and its plus factor
\begin{equation*}
  G^+(p)=
  \exp\!\left\{
    \frac{1}{2\pi\ii}
    \int_{-\ii\infty}^{\ii\infty}
    \frac{\log G(z)}{z-p}\,\dd z
  \right\},
  \qquad \operatorname{Re}p<0.
\end{equation*}
Here, $z$ is the complex integration variable, $\ii^2=-1$, and the contour
is directed upward along the imaginary axis. The gamma-function factor is
\begin{equation*}
  K^+(p)=\frac{\Gamma(1-p)}{\Gamma(\tfrac12-p)}.
\end{equation*}
Here, $\Gamma$ is the Euler gamma function. For real $p<0$, $G^+(p)$
can equivalently be evaluated using the real integral
\begin{equation*}
  \log G^+(p)
  =-\frac{p}{\pi}\int_0^\infty
    \frac{\log G(\ii t)}{t^2+p^2}\,\dd t.
\end{equation*}
Here, $t$ is a real integration variable.

\subsection{Corner with a process zone}

For a process zone in material~$i$, the kernel of the functional equation
in Section~\ref{sec:derivation} is
\begin{equation*}
  G_i(p)=\frac{D_i(p)\cos(\pi p)}{D(p)\sin(\pi p)}.
\end{equation*}
The corresponding characteristic determinants are
\begin{equation}
\label{eq:app-post-zone-determinants}
\begin{aligned}
  D_1(p)={}&2e(1+\kappa_2)
  [\cos(2p\alpha)-\cos2\alpha]
  \{\sin^2[p(\pi-\alpha)]-p^2\sin^2\alpha\}\\
  &+(1+\kappa_1)
  [\sin(2p\alpha)+p\sin2\alpha]
  [\sin\{2p(\pi-\alpha)\}-p\sin2\alpha],\\[3pt]
  D_2(p)={}&2(1+\kappa_1)
  [\cos\{2p(\pi-\alpha)\}-\cos2\alpha]
  [\sin^2(p\alpha)-p^2\sin^2\alpha]\\
  &+e(1+\kappa_2)
  [\sin\{2p(\pi-\alpha)\}-p\sin2\alpha]
  [\sin(2p\alpha)+p\sin2\alpha].
\end{aligned}
\end{equation}
% The author-confirmed signs in D_1 are preserved. Their provenance is
% recorded in EDITORIAL_LOG.md rather than in the clean manuscript.
The exponent $\lambda_i$ for a configuration with a zone in
material~$i$ is determined by
\begin{equation}
\label{eq:app-post-zone-root}
  D_i(-1-\lambda_i)=0,
  \qquad -1<\lambda_i<0,
  \qquad i=1,2.
\end{equation}
These roots define the $\lambda_1$ and $\lambda_2$ curves in
\SingularityExponentFigure. They characterize the corner singularity in the
presence of a process zone but are not used in the formulas for $d_i$,
$\delta_i$, and $J_i$; those formulas use $\lambda$ from
Eq.~\eqref{eq:app-D-root}.

\end{appendices}

% The following declarations are intentionally left as visible submission-preflight
% placeholders until all authors confirm the exact wording.
\section*{Statements and Declarations}
\textbf{Funding.} [To be confirmed before submission.]

\textbf{Competing interests.} [To be confirmed before submission.]

\textbf{Author contributions.} [To be confirmed by all authors before submission.]

\textbf{Data availability.} The numerical results reported in this study are obtained from the analytical relations and material parameters given in the article. Supporting calculation data are available from the corresponding author on reasonable request.

\bibliography{references}

\end{document}
